\documentclass[10pt,a4paper]{amsart}
\usepackage[margin=19mm]{geometry}
\usepackage{amsmath,amssymb,mathtools}
\usepackage[scr=rsfs,cal=euler]{mathalfa}
\usepackage{xfrac}
\usepackage[unicode,hidelinks,bookmarks=false]{hyperref}
\newcommand{\ie}{i.e.\ }
\newcommand{\cf}{cf.\ }
\newcommand{\resp}{respectively\ }
\newcommand{\Subsection}{\subsection}
\providecommand{\nobreakdash}{\nobreak\hbox{-}}
\setlength{\emergencystretch}{2em}
\multlinegap0pt
\usepackage{float}
\usepackage{amssymb}
\usepackage{enumerate}
\usepackage{subfig}
\usepackage{multirow}
\usepackage{xcolor}
\newtheorem{lem}{Lemma}
\title{An identity relating $n$-nacci numbers, partitions, and products of binomial coefficients}
\author{Du\v{s}an Dragutinovi\'c}
\date{Source: arXiv:2601.17401v1 (CC BY 4.0)}
\begin{document}
\maketitle

\noindent\emph{Source and licence.} An identity relating $n$-nacci numbers, partitions, and products of binomial coefficients. Version arXiv:2601.17401v1 (CC BY 4.0), distributed by arXiv under the Creative Commons Attribution 4.0 International licence, http://creativecommons.org/licenses/by/4.0/, as recorded in the arXiv licence metadata for this exact version and shown on the abstract page https://arxiv.org/abs/2601.17401v1. Full licence notice: \texttt{licenses/arxiv-2601.17401-CC-BY-4.0.txt}.\par\smallskip

\noindent\textit{Editorial scope.} Counted unit: the article's lem lem:grouping (source0.tex lines 491--494). The article's complete original proof of it is delivered with it (source0.tex lines 496--512, one \texttt{proof} environment of 17 lines). No mathematics is written, repaired or re-derived: the statement and the proof are the article's own lines, in the article's own notation, and no retained line is substituted or re-flowed.  Retained uncounted as the source-local context the proof needs: lines 237--240; lines 478--482.  Nothing was omitted from any retained span, so the package carries no omission record.  No retained line contains a citation, so the wrapper carries no bibliography and no bibliography entry.  No retained line contains a figure environment, an image-inclusion command or drawing code, so no image is delivered and the package declares no asset dependency.  Editorial formatting only, in addition: an AMS \texttt{amsart} preamble that restates the article's own declarations and macros used by the retained lines, namely the environments \texttt{lem} on separate counters, together with the standard packages those lines need.  The retained environments print in this document as Lemma 1 in order of appearance, whereas the article prints the same items with the numbering of its own full document.  The complete native source of the article as distributed in the arXiv e-print for this exact version is bundled unchanged as \texttt{sources/193312/source0.tex}.
\begin{equation}
  \nu =      (\underbrace{1, 2, \ldots, f}_{f}, \underbrace{*, \ldots, *, f}_{\delta_n}, \underbrace{*, \ldots, *, f + \delta_n}_{\delta_{n-1}}, \underbrace{*, \ldots, *, f + \sum_{i = n-1}^n\delta_i}_{\delta_{n-2}}, \ldots,  \underbrace{*, \ldots, *, f + \sum_{i = 2}^n \delta_i}_{\delta_1}),  
\label{eqn:nu_det_by_delta}
\end{equation}

\begin{lem}
Let $\delta = \delta_{1}\ldots\delta_m \vdash g$ and $\delta' = \delta_{1}'\ldots\delta_n' \vdash g$ be two partitions of $g$, and let $h \geq \delta_1$ and $f \leq \min(\delta_m, \delta_n')$ be positive integers. 
Then ${\delta} \leq_{ft} {\delta}'$ if and only if ${\Tilde{\delta}} \leq_{ft} {\Tilde{\delta'}}$, where $\Tilde{\delta} = h\delta_{1}\ldots\delta_m f$ and $\Tilde{\delta'} = h\delta_{1}'\ldots\delta_n'f$. 
\label{lemma:extension_delta}
\end{lem} 

\begin{lem}
Let $\delta = \delta_1\ldots\delta_{n + 1} \vdash g$ be a partition. Assume that $\delta' = \delta_1'\ldots\delta_n' \vdash g$ is a partition obtained by grouping two parts of $\delta$, i.e., such that multisets $\{\delta_1', \ldots, \delta_n'\}$ and $\{\delta_i + \delta_j\} \cup \{\delta_k: k \neq i, j\}$ coincide for some $i \neq j$. Then ${\delta'}\leq_{ft} {\delta}$.
\label{lem:grouping}
\end{lem}

\begin{proof} By Lemma \ref{lemma:extension_delta} and induction on $g$ and $n$, we may assume that $\delta_1' = \delta_1 + \delta_{n+1}$ and $\delta_i = \delta_i'$, for $2\leq i \leq n$. We need to show that for any $\nu'$ such that $\delta' = \delta'_{\nu'}$, there is a $\nu$ such that $\delta = \delta_{\nu}$ such that $\nu' \leq \nu$, i.e., that 
$\nu'(j) \leq \nu(j)$ for all $1\leq j \leq g$. 

For any final type $\Tilde{\nu}$, the following inequality holds: $\Tilde{\nu}(a + b) \leq \Tilde{\nu}(a) + b$, where $a, b \geq 0$; this follows directly from the definition. By \eqref{eqn:nu_det_by_delta}, for any $\nu'$ such that $\delta' = \delta'_{\nu'}$ and $\nu$ such that $\delta = \delta_{\nu}$, we thus have 
\begin{equation}
\nu'(\sum_{i = j}^{n}\delta_i) = \nu'(\sum_{i = j}^n\delta_i') = \sum_{i = j + 1}^{n}\delta_i = \nu(\delta_{n + 1} + \sum_{i = j}^{n}\delta_i) - \delta_{n + 1} \leq \nu(\sum_{i = j}^{n}\delta_i) = \nu(\sum_{i = j}^{n}\delta_i') \text{ for any }2\leq j \leq n.
\label{eqn:ft_comparison_1}
\end{equation}
Similarly, 
\begin{equation}
\nu'(\sum_{i = 1}^{n + 1}\delta_i) = \nu'(\sum_{i = 1}^n\delta_i') = \sum_{i = 2}^n\delta_i' = \sum_{i = 2}^n\delta_i \leq \delta_{n + 1} + \sum_{i = 2}^n\delta_i  = \nu(\sum_{i = 1}^{n + 1}\delta_i) = \nu(\sum_{i = 1}^{n}\delta_i').
\label{eqn:ft_comparison_2}
\end{equation}
We conclude that for any $\nu$ (resp.~$\nu'$) such that $\delta = \delta_{\nu}$ (resp.~$\delta' = \delta'_{\nu'}$), inequalities \eqref{eqn:ft_comparison_1} and~\eqref{eqn:ft_comparison_2} hold. Furthermore, all such $\nu\vdash g$ (resp.~$\nu'\vdash g$) can be any final types whose values appearing on the right-hand (resp.~left-hand) side of \eqref{eqn:ft_comparison_1} and \eqref{eqn:ft_comparison_2} are prescribed. 
Once again, recall that $\sum_{i = 1}^{n+1}\delta_i = \sum_{i = 1}^n\delta'_i$ and $\sum_{i = j}^n\delta_i = \sum_{i = j}^n\delta'_i$ for any $2 \leq j \leq n$.
Therefore, given any $\nu' \vdash g$ such that $\delta' = \delta'_{\nu'}$, we can always find $\nu \vdash g$ such that $\delta = \delta_{\nu}$ for which~$\nu' \leq \nu$.
\end{proof}

\end{document}
